\section{国内外研究现状}
\label{sec:research-status}

围绕“弱连接交往”与“大学生价值观”两个核心主题，本章按国内与国外两条脉络梳理既有研究成果。每一部分均先回顾弱连接交往的研究，重点说明弱关系如何传递信息、数字技术如何改变联系门槛以及当代青年在浅层社交中的真实体验；再回顾大学生价值观的研究，依次归纳青年价值观的现实特点、背后的影响因素以及高校的引导办法。最后综合评述中外研究的进展与不足，明确本研究的切入空间。

\subsection{国内研究现状}
\label{sec:domestic-research}

国内研究扎根于中国从传统熟人社会向现代流动社会转型的现实背景，并随着移动互联网普及和高校思政工作推进不断深入。整体来看，学者既关注弱关系在资源获取与校园“搭子”交往中的新特点，也持续追踪当代大学生价值观的变化趋势及其教育引导路径。

\subsubsection{对弱连接交往的研究}

国内关于弱连接交往的研究，大致沿着从传统熟人社会的关系结构分析，向数字媒介连接方式考察，再向当代大学生“搭子”与轻社交实践延伸的轨迹推进。

\paragraph{强弱关系与资源获取研究}
理解中国社会的弱连接，首先要回到本土人际关系的转型背景。费孝通在《乡土中国》中提出“差序格局”，指出传统乡土社会以血缘和地缘为纽带，人们主要依靠亲戚、邻里等熟人强关系获取生活资源与安全感，圈子内部情感深厚、人情往来密集，但对外较为封闭\cite{Fei2012}。随着市场化改革与人口跨区域流动加快，封闭的熟人圈子逐渐被打破。边燕杰较早通过对城市居民社会资本的调查发现，在重视人情信任的社会底色之上，跨越熟人圈子的弱关系开始发挥越来越重要的作用\cite{ZSHK200403011}；他在后来的关系社会学研究中进一步系统阐明，强关系主要承担信任担保功能，而弱关系擅长在不同群体之间牵线搭桥，能够帮助个体获取熟人圈子之外的新鲜信息与跨阶层职业机会\cite{Bian2021}。在具体衡量标准上，林聚任在《社会网络分析》中把互动频率、情感深度、亲密程度和互惠内容作为区分关系强弱的四项核心指标，从经验层面证实弱连接的独特价值在于打破小圈子的信息闭塞，成为沟通不同社会群体的桥梁\cite{Lin2009}。

\paragraph{数字媒介与交往方式研究}
智能手机和算法推荐的普及，改变了人与人建立联系的方式，也成为国内传播学与社会哲学关注的重点。彭兰分析指出，手机带来的“全时在线”状态模糊了工作、社交与私人生活的界限，青年在微信等熟人平台上承受着较重的人情维护和形象管理压力，因而产生收紧朋友圈权限甚至阶段性断联的“反连接”诉求\cite{GJXW201902002}；她随后进一步提出“液态连接”来概括数字化人际关系的新形态，指出如今的交往不再依赖稳固的熟人外壳，而是跟着具体场景和即时需求快速聚合，事情办完或场景转换便自然散去\cite{Peng2021}。蓝江从社会哲学角度将这一现象概括为“微粒化社会”，认为数字媒介拆解了家庭、邻里和单位等厚重的实体联系，把具体的人拆散成依靠兴趣标签临时聚合的独立个体，使人际往来逐渐失去共同生活积累下来的约束力\cite{Lan2022}。胡翼青和王敏则将这一变化归纳为日常交往的“轻量化转向”，指出年轻人在社交中越来越倾向于卸下沉重的人情负担，追求门槛低、维系省力、随时可以无痛退出的互动体验\cite{Hu2019}。从技术支撑来看，喻国明和曲慧在探讨算法推送机制时指出，智能算法通过对用户日常行为数据的抓取与分析，可以绕过现实熟人关系直接匹配细分兴趣节点，从底层重塑了弱连接的生成途径与信息流动方式\cite{XJSF202001015}。进一步针对青年网络圈层文化，蔡骐研究发现，靠算法和共同爱好聚起来的兴趣群体往往内部容易产生强烈共鸣，对外却筑起较高的沟通壁垒，反而加大了不同圈层之间相互理解的难度\cite{Cai2018}。

\paragraph{青年轻社交与搭子交往研究}
线上流动连接延伸到大学校园，催生了饭搭子、自习搭子、运动搭子、拼单搭子等广泛流行的“轻社交”现象。廉思基于全国性青年调查分析指出，“搭子社交”反映出青年从传统班级、社团等固定集体中逐渐抽离后，借助细分领域的简单约定建立起浅层陪伴关系；大学生试图在获得即时陪伴与减少人情责任之间寻找平衡，但长期依赖这种按功能拆分的交往，潜藏着削弱深度沟通意愿、降低共情能力和回避矛盾磨合的隐患\cite{Lian2024a}。结合升学与就业压力，孙寿涛和张晓芳对Z世代大学生的轻社交行为展开分析，指出这种选择并非单纯冷漠，而是年轻人在学业竞争激烈、未来充满不确定性的背景下，为节省精力、减少人际摩擦和精神内耗而采取的自我保护策略\cite{Sun2023}。这种心理防御在日常相处中表现为对人际距离的精细把控：刘航分析指出，搭子社交之所以在青年群体中迅速流行，是因为它精准迎合了当代大学生对清晰个人边界、低风险交流和实用互助的需求\cite{LiuHang2023}；陈姣和牛菲则提醒，如果过度依赖此类浅层关系，容易把原本带有情感温度的同学往来简化为按需调用、即时算清的实用交易\cite{ChenJiao2024}。邢海燕和邸涵针对大城市青年日常时空边界的调查也印证了这一趋势，表明当代青年习惯主动划清私人生活与公共社交的界线，排斥未经打招呼的熟人介入，借由高度可控的浅层接触来保持个人生活的自主与轻松\cite{Xing2022}。

\subsubsection{对大学生价值观的研究}

国内学界围绕大学生价值观的研究，紧扣改革开放以来的社会变迁与高校立德树人实践，形成了涵盖现状特点追踪、成因剖析与思政教育引导三个环节的完整脉络。

\paragraph{中国大学生价值观现状研究}
通过长期大样本调查，学者们较清晰地勾勒出中国大学生价值观的演变趋势。黄希庭和郑涌较早跳出照搬西方量表的局限，立足中国文化与社会现实，从人生目标、职业追求、人际伦理和公共责任等方面建立起本土青年价值观测量体系\cite{Huang2005}。金盛华和辛志勇在梳理中国人价值观发展趋势时同样指出，改革开放以来，青年价值观逐步从单一集体本位向兼顾个人发展与社会贡献的多元结构演变\cite{BJSF200303008}。近年来，廉思依托全国大样本调研将当代大学生的价值面貌概括为“宏微双轨”：在国家发展大局、制度认同和民族复兴等宏观立场上态度鲜明且高度稳固，但落到考研、求职、婚恋和日常相处等具体生活选择时，受竞争加剧与回报分化影响，价值取向明显变得务实且分化加剧\cite{Lian2024b}。沈壮海和王芸婷基于全国大学生思想政治状况滚动调查也证实，当代青年爱国底色鲜明，同时在学业深造、就业压力和人际适应上承受着突出的现实焦虑\cite{Shen2022}。冯刚和曾耀荣在综述近十年核心价值观培育研究时指出，正是这种宏观认同稳固与微观困惑并存的局面，推动学界从早期的宏观理论宣讲转向如何贴近大学生日常生活、回应具体关切的深入探索\cite{Feng2022}。在人际信任与亲密关系方面，辛自强和周正通过跨二十年的追踪对比分析发现，伴随人口高频流动与陌生人社会扩大，中国大学生对一般他人的信任度呈明显下降趋势\cite{XLXD201203004}；辛自强在后续关于青年社会心态的研究中进一步指出，青年在跨圈层交往中的戒备心理和防御边界持续收紧\cite{Xin2020}。风笑天针对青年婚恋观变迁的专项研究也相互佐证，表明当代大学生愈发看重精神独立、生活节奏自主和边界自由，对需要高度情感投入和长期责任捆绑的传统亲密关系态度更趋谨慎\cite{FengXiaotian2018}。

\paragraph{大学生价值观影响因素研究}
现有文献主要从社会结构转型和数字媒介环境两条线索，解释大学生价值观发生上述变化的原因。在社会结构层面，阎云翔在《中国社会的个体化》中提出“制度化个体主义”，指出市场化进程和用工制度改革促使青年从大家庭与传统单位的庇护中抽离出来，必须独自面对学业竞争与市场风险，因此个人权利、自主选择与自我负责迅速成为青年做判断时的核心标准\cite{Yan2012}。从代际成长环境看，周晓虹围绕“文化反哺”与代际转型指出，熟悉互联网的年轻一代在信息获取与文化消费上拥有更强的话语权，改变了传统的长幼单向传授格局\cite{Zhou2020}；李春玲从代际社会学视角刻画了中国“Z世代”的群体特征，说明成长于物质丰裕与互联网普及期的大学生普遍视野开阔，崇尚平等商量、反感居高临下的说教、尊重个体生活差异\cite{Li2021}。然而，陆士桢和蔡梦雅在剖析数字时代青年价值观新特征时指出，当独立的个体缺少稳固集体托底又面对激烈就业竞争时，不少青年在日常交友、社团活动和择业中更容易滑向看重即时回报与得失算计的实用主义心态\cite{Lu2022}。在媒介环境层面，胡正荣和李继东指出，移动互联时代去中心化的传播格局削弱了传统权威自上而下传递价值观的影响力\cite{Hu2017}。隋岩和曹飞从群体传播视角深入分析发现，离开实体班团集体的青年依托网络弱连接形成各类兴趣圈层，圈层内部高频的专属黑话互动与密集情绪感染，会在潜移默化中改写青年的是非标准与价值排序\cite{Sui2021}。朱哲和刘伟进一步探讨了网络圈层化对大学生核心价值观认同的冲击，认为商业平台算法制造的“信息茧房”助长了消费主义和泛娱乐化风气的扩散\cite{Zhu2021}；崔聪和张励仁则结合网络民粹思潮分析指出，缺乏事实核查和历史背景支撑的碎片化、情绪化信息流，容易让价值观尚未完全定型的部分学生陷入非黑即白、简单贴标签的偏激判断\cite{Cui2023}。

\paragraph{大学生价值观的思想政治教育引导研究}
面对个体化趋势与网络浅层社交带来的新挑战，国内思政学界在继承经典原理的基础上，探索了更加精准、贴近日常的育人路径。张耀灿等在《现代思想政治教育学》中阐明，主流价值观转化为学生内心的真实信念，需要经历“认知—认同—内化—外化”的渐进过程，唯有触动青年的现实交往体验与情感共鸣，外在规范才能真正落地生根\cite{Zhang2006}。沈壮海从文化视野出发，主张跳出单向说教，通过把宏大叙事转化为贴近青年生命的日常叙事、以润物无声的文化氛围回应学生的精神需求，打通价值认同的现实通道\cite{Shen2013}。刘建军则围绕社会主义核心价值观“日常化、具体化、形象化、生活化”展开论述，强调要把抽象原则转化为融入课堂、宿舍起居和同伴交往的具体行为准则，引导学生在日常互动中践行诚信与友善\cite{Liu2016}。在制度层面，教育部等部门先后印发《高校思想政治工作质量提升工程实施纲要》\cite{MOE2017_Gangyao}、《关于加快构建高校思想政治工作体系的意见》\cite{MOE2020_Gangshen}以及《全面推进“大思政课”建设的工作方案》\cite{MOE2022_Fangan}，对“十大育人”体系与社会大课堂融通作出系统部署。顺应这一导向，近年研究重点探索如何走进青年的网络弱连接圈层：骆郁廷和郭莉探讨了数字赋能“精准思政”的内在机理，主张借助大数据和智能技术动态感知学生在线上的社交轨迹与隐性困惑，实现分众化精准供给\cite{Luo2022}；戴艳军和孙晓晖聚焦话语转化，提出把严肃抽象的“文件话语”“学术话语”转化为贴近青年交流习惯的“生活话语”“互动话语”\cite{Dai2021}；杨晓慧则围绕数字时代高校网络思政的分众化引导提出，要将线上弱连接圈层的柔性融入与线下校园共同体（如一站式学生社区、科研团队和社会实践）的情感温度重构结合起来，在虚实协同中重新培植大学生的集体归属感与社会责任感\cite{Yang2023}。

\subsection{国外研究现状}
\label{sec:international-research}

国外对弱连接交往与青年价值观的研究起步较早，横跨社会学、心理学、传播学、计算社会科学与教育学等多个领域，既积累了经典的理论概念，也开展了大量基于真实数据和课堂实验的实证检验。

\subsubsection{对弱连接交往的研究}

国外关于弱连接交往的研究，重点围绕弱关系如何带来信息与机会、互联网如何降低连接成本，以及浅层交往给现代人带来怎样的情绪体验与关系困境三个问题展开。

\paragraph{弱连接与机会获取}
西方古典社会学很早就注意到现代人际关系由紧密走向松散的趋势。滕尼斯（Tönnies）在《共同体与社会》中对比了基于血缘、邻里默契的传统“共同体”和基于理性算计、明确契约的现代“社会”，指出现代化使人际纽带从全面照顾转向特定功能合作\cite{Tonnies2019}。齐美尔（Simmel）在分析现代大都市生活时指出，分工细化让个体同时处于多个松散交叠的社交圈子里，这既带来了前所未有的行动自由，也拉开了人与人之间的心理距离，催生出都市人理智算计与冷漠防御的心态\cite{Simmel2002}。真正奠定该领域理论基础的是美国社会学家格兰诺维特（Granovetter），他在经典论文《弱关系的力量》中从互动频率、情感深度、亲密程度和互惠交换四个方面衡量关系强弱，证明联系紧密的强关系大多圈在背景相似的熟人内部，信息高度重复；而联系较少、感情较浅的弱关系能跨越小圈子充当“信息桥”，为人们带来跨行业的求职信息与新思路\cite{Granovetter1973}。顺着这一思路，伯特（Burt）提出“结构洞”理论，说明处于两个互不相连群体之间的中间人，能够优先获取不同圈子的新信息并掌握协调资源的主动权\cite{Burt1992}。帕特南（Putnam）则把社会资本区分为凝聚熟人内部的“粘合型社会资本”和连接不同背景人群的“桥接型社会资本”，指出后者对于打破群体偏见、拓宽公共视野和培育广泛社会信任具有关键作用\cite{Putnam2000}。近五年，计算社会科学用超大规模数据检验了这些经典判断：拉杰库马尔等（Rajkumar et al.）在职业社交平台上对2000万名用户开展长达五年的随机对照实验，证实中等强度的弱连接最能促进青年跨行业求职与职业流动\cite{Rajkumar2022}；切蒂等（Chetty et al.）追踪分析7200万名社交媒体用户（涵盖全美绝大多数高校学生）后发现，跨越不同收入阶层的弱连接（即“经济连通性”），比小群体内部团结或参加社团更能提升青年的学业抱负与向上流动机会\cite{Chetty2022}。

\paragraph{数字媒介与弱连接的形成}
互联网与社交媒体大幅压低了建立和维护弱连接的门槛。海桑思韦特（Haythornthwaite）提出“潜在关系”概念，指出大学校园或大型机构里原本就存在大量因同校、同城或同群而具备联系可能、却尚未实际往来的人，数字媒介极大降低了寻找和破冰的成本，让大学生可以按需把这些沉睡的联系随时变成现实弱连接\cite{Haythornthwaite2005}。埃里森等（Ellison et al.）对美国公立大学本科生展开追踪调查，发现大学生在社交媒体上维系广泛的校园“泛泛之交”，能有效积累桥接型社会资本，尤其有助于性格内向或有社交焦虑的新生适应异地求学生活、增强校园归属感\cite{Ellison2007}。雷尼和威尔曼（Rainie \& Wellman）将这种数字时代的社交形态概括为“网络化个体主义”，说明年轻人正在从地域和家族等固定捆绑的社群中走出来，转而以自己为中心，根据不同需要灵活编织跨越空间的松散联系网\cite{Rainie2012}。拜耳等（Bayer et al.）的系统综述进一步揭示了支撑这种网络的日常机制，指出社交平台提供的各类功能设计——比如阅后即焚、分组可见、兴趣群聊和轻量点赞——让年轻人能够精细地控制社交距离、在线状态和自我暴露的深浅\cite{Bayer2020}。

\paragraph{弱连接的交往体验与关系维系}
关于弱连接给现代人带来的心理体验与长远影响，国外行为心理学和社会批判理论给出了截然不同却彼此互补的观察。在积极一面，桑德斯特伦和邓恩（Sandstrom \& Dunn）通过校园追踪实验发现，大学生在日常生活中与不太熟悉的同学、社团外围成员或咖啡店员进行简短友善的闲聊，就能显著提升当天的愉悦感和归属感\cite{Sandstrom2014}；阿蒂尔等（Atir et al.）针对陌生人交谈开展系列行为实验进一步证明，人们普遍低估了与陌生人或弱关系对象聊天所能带来的信息收获与好心情，只要迈出搭话的第一步，往往能打破固有局限、获得意外启发\cite{doi:10.1073/pnas.2206992119}。在隐忧一面，社会理论家深入剖析了浅层连接过度蔓延后对深度关系的侵蚀。吉登斯（Giddens）指出，现代社会让人际关系脱离了传统宗族和地域的外部约束，演化成全靠双方当下心理满意来维系、随时可以单方面退出的“纯粹关系”\cite{Giddens1998}。鲍曼（Bauman）在《液态之爱》中更为尖锐地批判，消费主义逻辑已经渗入日常交友，诱导现代人像挑选快消品一样建立随时可开启、随时可无痛终止的“轻盈连接”，让人一边贪恋低成本的陪伴，一边本能地逃避需要磨合与承担责任的深度关系\cite{Bauman2003}。罗萨（Rosa）从时间节奏的角度提出“社会加速”理论，认为科技、社会变化与生活步调不断提速，迫使年轻人在碎片时间里塞满大量浅层社交信息，挤干了静心倾听和建立深层情感共鸣所需的时间，使高频联系的背后弥漫着空虚感\cite{Rosa2015}。这一判断得到了实证研究的有力印证：特克尔（Turkle）通过长期访谈北美高校学生发现，大学生越来越习惯用可以反复修改、零情绪风险的文字消息替代面对面交流，长此以往导致察言观色的共情能力、深度倾听耐心以及包容不同意见的雅量明显退化\cite{Turkle2015}；霍尔（Hall）的实证测算更直观地揭示了时间挤占问题——从泛泛之交变成普通朋友平均需要约50小时的真实相处，变成亲密挚友则需要200小时以上的高质量陪伴，大学生把大量闲暇耗在刷动态和点赞等浅层连接上，客观上挤占了培养真挚友谊的稀缺时间\cite{Hall2019}。

\subsubsection{对大学生及青年价值观的研究}

国外学者结合青年成长阶段与数字时代特征，围绕大学生价值观的测量与特点、环境与媒介成因以及学校教育的引导机制，开展了大量跨文化调查与课堂实验研究。

\paragraph{大学生及青年价值观现状与特点研究}
西方发展心理学和跨文化研究系统刻画了当代大学生的价值特征。阿内特（Arnett）提出“初显期成人”（18—25岁）理论，指出高等教育延长使大学生脱离了未成年时的家庭依附，又尚未步入稳定婚姻和终身职业，处于探索自我、流动性高且高度关注自身发展的独特时期，因此格外看重自主决策和多元体验\cite{Arnett2000}；他在后续修订著作中进一步证实，这一阶段强调情感自主与个性探索的价值倾向在当代各国青年中愈发普遍\cite{Arnett2023}。在价值观分类测量方面，早期心理学家斯普朗格（Spranger）从人格整体出发，将人类主导价值划分为理论型、经济型、审美型、社会型、政治型和宗教型六类\cite{Spranger2017}。在此基础上，施瓦茨（Schwartz）构建了极具国际影响力的“人类基本价值观环形模型”，用“关爱他人还是追求个人成功”“乐于求变还是看重传统稳定”两组对照维度划分出十类基本价值观\cite{Schwartz1992}；施瓦茨等（Schwartz et al.）在2012年进一步将模型细化为十九类，并基于跨文化样本证实，大学阶段青年在自主思考、乐于尝试和普世包容上的权重显著高于中老年群体，而在遵守传统与顺从规范上得分偏低\cite{Schwartz2012}。政治社会学家英格尔哈特（Inglehart）的长期追踪也印证了这一趋势，认为成长于物质相对充裕年代的年轻一代，价值重心正从关注温饱与安全转向强调自我表达和生活质量的“后物质主义”\cite{Inglehart2018}。然而，2010年代智能手机和社交媒体普及之后，情况出现了新的复杂变化。特温格（Twenge）基于全美数百万大中学生跨世代追踪数据发现，当代“手机世代”呈现出“高包容与高脆弱”并存的特征：他们在平等、多元文化等议题上态度开放，但线下同伴聚会、打工实践和独立出行的比例大幅下滑，孤独感与抑郁、焦虑等心理脆弱指标显著上升\cite{Twenge2017}。卢基安诺夫和海特（Lukianoff \& Haidt）指出，这促使不少大学生把“避免心理不适”奉为高于一切的准则（情感安全主义），习惯把日常观点分歧放大为严重伤害，面对不同意见时倾向于回避或拉黑\cite{Lukianoff2018}；海特（Haidt）在2024年的最新著作中进一步追溯其根源，论证青年成长方式从线下自由玩耍、在同伴小摩擦中学会协商，全面转向手机屏幕与虚拟弱连接，剥夺了年轻人在真实交往中锻炼抗挫折能力的机会\cite{Haidt2024}。

\paragraph{价值观影响因素研究}
国外学界从成长环境、代际经历、社会交往规则以及平台微观机制四个层次剖析青年价值观的成因。罗克奇（Rokeach）较早将价值观区分为追求理想人生状态的“终极价值观”和指导日常做法的“工具价值观”，说明社会制度与文化规范主要通过调整这两类价值的轻重排序来影响个人判断\cite{Rokeach1973}。在环境层级上，布朗芬布伦纳（Bronfenbrenner）提出“人类发展生态学模型”，强调青年价值观在家庭与同伴圈、学校、社区以及宏观社会文化的层层互动中形成\cite{Bronfenbrenner1979}；顺应移动互联时代，纳瓦罗和图吉（Navarro \& Tudge）提出“新生态理论”作出修正，将青年日常所处的微观环境拆分为面对面的现实空间和跨时空的网络空间，指出社交媒体上的高频弱连接互动已经演化为独立塑造大学生价值认知与自我认同的成长空间\cite{Navarro2023}。在时间维度上，曼海姆（Mannheim）的代际理论指出，17至25岁是青年脱离原生家庭习惯、初次独立接触社会并形成稳定政治与社会价值观的“关键敏感期”，这一时期共同经历的技术变革与社会事件会沉淀为整代人的价值底色\cite{Mannheim1952}。在深层交往规则层面，哈贝马斯（Habermas）在《交往行为理论》中批判了市场和官僚体制中的效率算计不断侵入日常生活的现象，指出这种扩张正在蚕食以真诚沟通和相互理解为核心的正常交往，为审视当代青年社交中的功利算计提供了深刻视角\cite{Habermas2004}。在此基础上，近年的实证研究查明了网络弱连接扭曲青年价值判断的三重微观机制：其一，桑斯坦（Sunstein）指出，算法个性化推送筑起的“回音室”不断削弱社会成员共享的公共经验基础\cite{Sunstein2017}，而贝尔等（Bail et al.）通过现场随机实验发现，在缺乏熟人信任托底的社交媒体弱连接网络中，单靠算法让用户接触对立阵营的陌生人观点，反而容易激发强烈的群体防御心理，加剧立场极端化\cite{Bail2018}；其二，布雷迪等（Brady et al.）对海量社交推文的大数据分析表明，在松散的弱连接传播网络中，带有激烈道德指责与愤怒情绪的言论比温和说理更容易获得点赞和转发，平台的可见性奖励不断诱导青年在公共议题上抛弃全面考量，采取更具攻击性的偏激立场\cite{Brady2021}；其三，韦尔杜因等（Verduyn et al.）的系统综述证实，由于线上弱连接网络充斥着泛泛之交精心修饰的“高光生活”，大学生在缺乏真实背景了解的情况下被动浏览这些动态，极易陷入无休止的向上社会比较、心理落差与“错失恐惧”，潜移默化地助长物质攀比并削弱自我价值感\cite{Verduyn2020}。

\paragraph{教育对价值观的塑造作用研究}
针对个体化与数字化带来的青年价值失范风险，西方教育学、道德心理学和媒介素养学者探索了多种具体干预路径。在集体生活与实践育人方面，涂尔干（Durkheim）在《道德教育》中指出，学校道德教育的关键在于培养纪律精神、关心集体的利他情怀和审慎理解规则的自主能力，强调唯有置身真实、有机的学校集体生活，才能克服社会原子化带来的利己倾向\cite{Durkheim2001}；杜威（Dewey）则提出“学校即社会”，批评脱离实际交往的抽象说教，论证真正有生命力的公共责任感只能生长于学生共同合作解决实际问题的日常交往之中\cite{Dewey2013}。在道德心理与情感培育方面，科尔伯格（Kohlberg）通过长期追踪提出道德认知发展的六阶段理论，证明引导青年超越功利交换迈向公正原则的关键，在于教师有意识地组织“道德两难讨论”，借助同伴之间的真实观点碰撞打破自我中心思维\cite{Kohlberg1981}；诺丁斯（Noddings）从“关怀伦理学”角度补充指出，化解现代人际冷漠不能只靠抽象推理，还要把教育重心落回具体的人际关怀，通过教师的榜样示范、真诚平等的对话、制度化的关怀实践以及对学生善意动机的肯定，培养青年体谅他人的共情能力与责任感\cite{Noddings2013}。在高校实证评估方面，梅休等（Mayhew et al.）整合了21世纪以来全美数千项高校追踪调查，证实大学之所以能显著提升学生的利他精神、公共责任感与包容度，最核心的推动力来自课堂之外高参与度的同伴真实合作、跨群体项目、住宿学院共同体生活以及深入社区的“服务性学习”，这些深度协作实践能够有效对冲浅层连接带来的疏离与功利心态\cite{Mayhew2016}。面对网络弱连接环境中的算法偏见与信息误导，崔等（Choi et al.）的调查表明，教师自身的网络实践经验与数字伦理意识，直接决定了学校数字公民课程能否切中学生的现实痛点\cite{Choi2018}；温伯格等（Wineburg et al.）近年在多所高校开展的课堂对照实验进一步证明，在通识课中系统教授专业事实核查员使用的“横向阅读”方法——即遇到陌生网络信息时主动跳出当前页面去核查信源背景与利益关联——能够显著提高大学生辨别误导信息的能力，培养审慎求真、不盲从的理性判断习惯\cite{Wineburg2022}。

\subsection{国内外研究综合述评}
\label{sec:research-review}

综观国内外既有研究，学界围绕弱连接交往演变与大学生价值观变迁已取得丰硕成果，为本研究奠定了扎实的理论与方法基础。中外研究一致证实，在社会个体化转型与移动互联网普及的共同推动下，以情感投入浅、边界清晰、退出成本低和按需匹配为特征的弱连接交往（包括校园“搭子”、轻社交以及线上跨圈层联系），已成为当代大学生日常互动的重要方式。既有文献通过大规模数据实验、长期追踪调查和深度访谈，既证实了弱连接在拓宽信息来源、助力职业流动、缓解新生适应焦虑和提供轻松情绪慰藉方面的积极作用，也指出了浅层连接过度蔓延背后潜藏的人际功利化、社会信任下降、共情能力退化、向上比较焦虑以及公共讨论极化等隐忧。总体而言，国外研究长于用大规模数据和课堂对照实验检验因果机制与干预效果；国内研究则紧密立足中国从熟人社会向流动社会转型的本土脉络，精准把握当代大学生价值观“宏微双轨”的现实特点，并依托“十大育人”体系、“精准思政”和“大思政课”探索价值引领路径。

然而，结合本研究的主题加以审视，既有文献在以下三个方面仍有进一步深化的空间。第一，在理论视角上，社会网络分析与思想政治教育价值引导仍相对割裂。社会学和传播学关于弱连接的研究多聚焦求职信息传递、社会资本积累或即时情绪获取的实用功能，往往把行动者的价值观念当成既定前提，抑或停留在宏观现代性批判层面；思想政治教育领域的网络研究则多集中于网络不良思潮辨析和话语转换技巧，较少回到马克思主义关于“人的本质是一切社会关系的总和”这一基本立场，系统剖析弱连接这种日常交往形式本身，如何通过潜移默化的相处规则一步步改变大学生的价值判断尺度，进而冲击马克思主义价值观与社会主义核心价值观的培育。第二，在作用过程上，缺少跨生活领域的整体整合。受学科分工限制，现有实证研究大多只截取青年生活的某一个切面（例如单独讨论“搭子社交”的心理防御，或单独检验社交媒体的向上比较），尚未把大学生的个人消费休闲、学业职业追求、同伴及亲属往来与社会公共参与作为一个相互牵动的整体加以考察，也未能充分揭示信息片面展示、低义务相处习惯、浅层满足替代与松散群体环境是如何共同作用于日常价值判断的。第三，在应对策略上，针对流动社交特点的系统方案尚显不足。既有引导措施容易走向两个极端：要么依赖传统行政班级与刚性集体活动去约束学生的浅层社交，因忽视青年合理的边界诉求而引发逆反；要么一味顺应算法推荐去迎合碎片化、轻量化偏好，无形中强化了工具算计与圈层封闭。学界仍缺乏一套既尊重当代青年交往特点、又能通过重构校园实体协作共同体和强化马克思主义方法论训练来化解价值偏差的系统应对路径。这正是本研究力图推进和突破的方向。